%! Function Composition and Decomposition — Unit 1, Lesson 1.3 — Homework
% THE LESSON'S GRADED PRACTICE and the last component in the packet — exactly TWO pages, started
% in the Homework Launch (the last 3 minutes of the period) and finished at home. Every context
% here is new: the notes used the Summit Cable Car; the AP Practice used a bus's fuel use and a
% wind farm; the homework uses a step counter and a shipping desk.
%   Page 1 — Part A (handoffs, 1–2) and Part B (building composites, 3–4).
%   Page 2 — Part C (chains in context, 5–6) and Part D (decomposition and the headline, 7–9).
% The diagnostic item is problem 2: one input, both orders, two different answers. Problem 9
% asks for the lesson's headline in writing. Problem 8 is the on-ramp to 1.4.
%
% PAGE LOCKSTEP with the other half of the pair: work blocks are byte-identical, prose answers
% are \writespace{H}{…} (fixed height both ways), and each \blank sits at the END of a line so
% its \ans cannot rewrap anything. The single explicit \newpage fixes the break in both files.
\documentclass[10pt]{article}
\usepackage{precalculus-article}
\usepackage{precalculus-boxes}
\newcommand{\TallMath}[1]{$\displaystyle #1\rule[-1.4em]{0pt}{3.2em}$}

\begin{document}

\pageheader{Unit 1, Lesson 1.3}{Homework}

\vspace{0.1in}

\boxguard[8]
\begin{remindbox}
\small
\textbf{This is your graded homework.} Start it in class; finish it on your own by the due
date on your cover sheet. Show your work, and on every composite \emph{write the handoff down}
before you use it --- the middle number is where the credit is.
\end{remindbox}

\vspace{0.1in}

\begin{headlinebox}{lilac}
\textbf{\color{plum}Part A --- Handoffs}
\end{headlinebox}

\begin{enumerate}[leftmargin=1.9em, itemsep=8pt]

\item \begin{minipage}[t]{0.40\linewidth}
      \vspace{0pt}
      The table gives every value of two functions $f$ and $g$.
      \end{minipage}\hfill
      \begin{minipage}[t]{0.56\linewidth}
      \vspace{0pt}
      \centering
      {\renewcommand{\arraystretch}{1.3}
      \begin{tabular}{|c|c|c|c|c|c|}
      \hline
      $x$ & 1 & 2 & 3 & 4 & 5 \\ \hline
      $f(x)$ & 4 & 1 & 5 & 2 & 3 \\ \hline
      $g(x)$ & 3 & 5 & 1 & 4 & 9 \\ \hline
      \end{tabular}}
      \end{minipage}

      \textbf{(a)} $f(g(1))$: \; handoff $= $ \blank{1.4cm}, \; so $f(g(1)) = $ \blank{1.6cm}

      \textbf{(b)} $g(f(1))$: \; handoff $= $ \blank{1.4cm}, \; so $g(f(1)) = $ \blank{1.6cm}

      \textbf{(c)} $f(g(5))$ has no value. In one sentence, say what goes wrong with the handoff.
      \writespace{0.9cm}{}

\item Let $f(x) = 4x$ and $g(x) = x - 5$. Run the input 8 through both orders.

\begin{work}
  f(g(8)) &= f(3) = 12 \\
  g(f(8)) &= g(32) = 27
\end{work}

      \textbf{(a)} $f(g(8))$: \; handoff $= $ \blank{1.4cm}, \; so $f(g(8)) = $ \blank{1.6cm}

      \textbf{(b)} $g(f(8))$: \; handoff $= $ \blank{1.4cm}, \; so $g(f(8)) = $ \blank{1.6cm}

      \textbf{(c)} Same two rules, same input. Are the two answers equal? \blank{2.2cm}

\end{enumerate}

\vspace{0.1in}

\begin{headlinebox}{lilac}
\textbf{\color{plum}Part B --- Building the composite}
\end{headlinebox}

\begin{enumerate}[leftmargin=1.9em, itemsep=8pt, start=3]

\item Let $f(x) = x^2 + 2$ and $g(x) = x - 3$. Build both composites and simplify.

\begin{work}
  f(g(x)) &= f(x - 3) = (x-3)^2 + 2 \\
          &= x^2 - 6x + 11
\end{work}

      \textbf{(a)} $(f \circ g)(x) = $ \blank{4.0cm} \qquad
      \textbf{(b)} $(g \circ f)(x) = $ \blank{3.2cm}

\item Let $f(x) = 5x + 1$ and $g(x) = 2x - 4$.

      \textbf{(a)} $(f \circ g)(x) = $ \blank{3.2cm} \qquad
      \textbf{(b)} $(g \circ f)(x) = $ \blank{3.0cm}

\end{enumerate}

\newpage

\begin{headlinebox}{lilac}
\textbf{\color{plum}Part C --- Chains in context}
\end{headlinebox}

\begin{enumerate}[leftmargin=1.9em, itemsep=8pt, start=5]

\item A step counter converts steps to distance: $K(s) = 0.0008s$ gives the kilometers walked in
      $s$ steps. A fitness app converts distance to energy: $C(k) = 60k$ gives the calories burned
      walking $k$ kilometers.

      Find $C(K(10{,}000))$. Name the handoff and its units on the way.
      \writespace{1.2cm}{}

\item A shipping desk converts a package's weight from ounces to pounds with
      $P(z) = \dfrac{z}{16}$, and prices a package by weight with $C(p) = 3 + 2p$ dollars for
      $p$ pounds. Explain why $C(P(48))$ is meaningful but $P(C(48))$ is not.
      \writespace{1.2cm}{}

\end{enumerate}

\vspace{0.1in}

\begin{headlinebox}{lilac}
\textbf{\color{plum}Part D --- Taking it apart, and the headline}
\end{headlinebox}

\begin{enumerate}[leftmargin=1.9em, itemsep=8pt, start=7]

\item Take each function apart into an inner rule and an outer rule.

      \textbf{(a)} $h(x) = (4x - 1)^5$: \quad inner $g(x) = $ \blank{2.4cm}, \quad
      outer $f(u) = $ \blank{2.2cm}

      \textbf{(b)} $k(x) = \dfrac{5}{2x - 9}$: \quad inner $g(x) = $ \blank{2.4cm}, \quad
      outer $f(u) = $ \blank{2.2cm}

\item Let $f(x) = 3x + 6$ and $g(x) = \dfrac{x - 6}{3}$.

\begin{work}
  f(g(x)) &= 3\!\left(\frac{x-6}{3}\right) + 6 = x \\
  g(f(x)) &= \frac{(3x+6) - 6}{3} = x
\end{work}

      $(f \circ g)(x) = $ \blank{2.0cm} \qquad $(g \circ f)(x) = $ \blank{2.0cm}

      What do these two functions do to each other?
      \writespace{0.9cm}{}

\item A classmate says $f(g(x))$ and $g(f(x))$ must be the same function because they are built
      from the same two rules. Using the word \emph{handoff}, write the sentence that settles it.
      \writespace{1.2cm}{}

\end{enumerate}

\vspace{0.08in}

\boxguard[8]
\begin{spiralbox}
\small
\textbf{Next lesson (1.4, Inverse Functions):} problem 8 found a pair of functions that undo
each other. Next lesson names that partner $f^{-1}$, and the test for it is a composition run
in \emph{both} directions: $f(f^{-1}(x)) = x$ \emph{and} $f^{-1}(f(x)) = x$.
\end{spiralbox}

\end{document}
